\chapter{On-Policy, Off-Policy, and the Sampling Problem}
\label{ch:on-off-policy}

Whether an algorithm is on-policy or off-policy is not a label to memorize per algorithm name. It is a consequence of a single question about the algorithm's update rule. This chapter derives the answer from first principles for every major algorithm family, then examines the practical implications for data efficiency.

%% ----------------------------------------------------------------
\section{The First-Principles Criterion}
\label{sec:criterion}

To classify any RL algorithm, ask one question:

\begin{keyidea}[The Sampling Criterion]
Does the update formula require the actions in the training data to have been sampled by the \emph{current} policy $\pi_\theta$?
\end{keyidea}

More precisely, check whether two chains are consistent:
\begin{itemize}
  \item \textbf{Sampling chain:} Which policy selected the actions in the data?
  \item \textbf{Optimization chain:} Which policy's parameters are being updated?
\end{itemize}

This yields three categories:
\begin{center}
\begin{tabular}{@{} l l @{}}
\toprule
\textbf{Category} & \textbf{Requirement on data actions} \\
\midrule
On-policy & $a \sim \pi_\theta(a \mid s)$ \quad (current policy) \\
Off-policy & $a \sim \mu(a \mid s)$ \quad (any behavior policy) \\
Near on-policy & $a \sim \pi_{\text{old}}(a \mid s)$, with $\pi_{\text{old}} \approx \pi_\theta$ enforced \\
\bottomrule
\end{tabular}
\end{center}

The real criterion is not ``whether there is a replay buffer'' but whether the action distribution in the data must match the distribution of the policy currently being optimized.

%% ----------------------------------------------------------------
\section{Why Policy Gradient Is On-Policy}
\label{sec:pg-on}

The policy gradient takes the form:
\begin{equation}
  \nabla_\theta J(\theta)
  = \mathbb{E}_{s,a \,\sim\, \pi_\theta}\!\bigl[
      \nabla_\theta \log\pi_\theta(a \mid s)\;\hat{A}^{\pi_\theta}(s,a)
    \bigr].
  \label{eq:pg-onpolicy}
\end{equation}

The expectation is taken over $s, a \sim \pi_\theta$. This is not a notational convenience; it is a mathematical requirement. The derivation uses the identity $\nabla_\theta \pi_\theta(a \mid s) = \pi_\theta(a \mid s)\,\nabla_\theta\log\pi_\theta(a \mid s)$, which only produces an unbiased gradient estimate when actions are drawn from $\pi_\theta$.

If actions came from an old or different policy $\mu \neq \pi_\theta$, the weighted log-likelihood $\log\pi_\theta(a \mid s)\cdot\hat{A}(s,a)$ is no longer an unbiased estimate of $\nabla_\theta J(\theta)$.

\begin{keyidea}[PG Is On-Policy by Derivation]
REINFORCE, A2C, and A3C are on-policy not because of their names, but because the gradient formula~\eqref{eq:pg-onpolicy} requires $a$ to come from the current $\pi_\theta$'s sampling distribution.
\end{keyidea}

%% ----------------------------------------------------------------
\section{Why Q-Learning Is Off-Policy}
\label{sec:ql-off}

The Q-learning target is:
\begin{equation}
  y = r + \gamma \max_{a'} Q_{\bar\theta}(s', a').
  \label{eq:ql-target}
\end{equation}

Crucially, this target does not depend on how action $a$ was selected. The transition $(s, a, r, s')$ tells us: ``someone executed $a$ in $s$ and observed $r, s'$.'' The target then asks: ``if the agent acts greedily from $s'$ onward, what is the expected return?'' The identity of ``someone'' is irrelevant---it could be the current policy, an old policy, a human operator, or a random controller.

\begin{keyidea}[Bellman Doesn't Care Who Sampled]
Q-learning's target involves $\max_{a'} Q(s', a')$, which is independent of the behavior policy that generated the data. This is why Q-learning can learn from a replay buffer of transitions collected by arbitrary policies---the fundamental reason it is off-policy.
\end{keyidea}

The same logic extends to DQN and Double DQN: the Bellman target depends on the \emph{target network's} value estimates, not on the policy that produced the stored transitions.

%% ----------------------------------------------------------------
\section{Why PPO Is Near On-Policy}
\label{sec:ppo-near}

PPO collects data with a recent snapshot of the policy $\pi_{\text{old}}$, then performs multiple gradient steps on the current policy $\pi_\theta$. This is not purely on-policy (the data was not sampled by $\pi_\theta$) nor purely off-policy (PPO cannot use arbitrary historical data).

PPO corrects for the distribution mismatch using an importance ratio:
\begin{equation}
  r_t(\theta) = \frac{\pi_\theta(a_t \mid s_t)}{\pi_{\text{old}}(a_t \mid s_t)},
\end{equation}
and then clips it to prevent large updates:
\begin{equation}
  \mathcal{L}_{\text{PPO}} = -\min\!\bigl(
    r_t(\theta)\,\hat{A}_t,\;
    \operatorname{clip}(r_t(\theta),\, 1{-}\epsilon,\, 1{+}\epsilon)\,\hat{A}_t
  \bigr).
  \label{eq:ppo-clip}
\end{equation}

The clipping enforces that $\pi_\theta$ and $\pi_{\text{old}}$ cannot diverge too much. Data from the most recent old policy is acceptable; data from a much older policy is not.

\begin{keyidea}[Near On-Policy = Bounded Staleness]
PPO allows limited reuse of recent data, but constrains distribution shift through ratio clipping. It is neither on-policy (data is slightly stale) nor off-policy (it cannot use arbitrary historical data). The precise term is \emph{near on-policy}.
\end{keyidea}

TRPO follows the same logic but enforces the constraint via a KL-divergence trust region rather than clipping.

%% ----------------------------------------------------------------
\section{Why DDPG, TD3, and SAC Are Off-Policy}
\label{sec:offpolicy-ac}

All three algorithms use a replay buffer, but the deeper reason for their off-policy nature is in their update rules.

\subsection{Critic Update}

The critic targets have the same structure as Q-learning---a Bellman backup that does not depend on the behavior policy:
\begin{align}
  y_{\text{DDPG}} &= r + \gamma\, Q_{\bar\phi}\!\bigl(s',\, \mu_{\bar\theta}(s')\bigr), \\
  y_{\text{TD3}}  &= r + \gamma\, \min_{i=1,2} Q_{\bar\phi_i}\!\bigl(s',\, \mu_{\bar\theta}(s') + \tilde\epsilon\bigr), \\
  y_{\text{SAC}}  &= r + \gamma\, \min_{i=1,2} Q_{\bar\phi_i}(s', \tilde{a}') - \alpha\log\pi_\theta(\tilde{a}' \mid s'),
  \quad \tilde{a}' \sim \pi_\theta(\cdot \mid s').
\end{align}

In each case, the action $a$ stored in the replay buffer can come from any previous version of the policy. The target is constructed from target networks and does not require $a \sim \pi_\theta$.

\subsection{Actor Update}

The actor optimizes \emph{through} the critic:
\begin{equation}
  \nabla_\theta J \approx \nabla_\theta\, Q_\phi\!\bigl(s,\, \mu_\theta(s)\bigr)
  \quad\text{(DDPG/TD3)},
  \qquad
  \nabla_\theta\, \mathbb{E}_{a\sim\pi_\theta}\!\bigl[Q_\phi(s,a) - \alpha\log\pi_\theta(a \mid s)\bigr]
  \quad\text{(SAC)}.
\end{equation}
The states $s$ are drawn from the replay buffer, but the actions used in the gradient are freshly computed by the current actor/policy. The actor does not need the stored action to match $\pi_\theta$---it generates its own actions at update time.

%% ----------------------------------------------------------------
\section{Quick Reference Table}
\label{sec:reference-table}

\begin{center}
\begin{tabular}{@{} l l l l @{}}
\toprule
\textbf{Algorithm} & \textbf{Data source} & \textbf{Requires $a \sim \pi_\theta$?} & \textbf{Verdict} \\
\midrule
REINFORCE   & Current policy       & Yes          & On-policy \\
A2C / A3C   & Current policy       & Yes          & On-policy \\
PPO         & Recent old policy    & Approximately & Near on-policy \\
TRPO        & Recent old policy    & Approximately & Near on-policy \\
Q-learning  & Any behavior policy  & No           & Off-policy \\
DQN         & Replay buffer        & No           & Off-policy \\
Double DQN  & Replay buffer        & No           & Off-policy \\
DDPG        & Replay buffer        & No           & Off-policy \\
TD3         & Replay buffer        & No           & Off-policy \\
SAC         & Replay buffer        & No           & Off-policy \\
\bottomrule
\end{tabular}
\end{center}

%% ----------------------------------------------------------------
\section{Data Efficiency: Theory vs.\ Engineering}
\label{sec:data-efficiency}

\subsection{Theoretical Sample Efficiency}

Off-policy methods can replay each transition many times, giving them a clear theoretical advantage in sample efficiency. A single $(s, a, r, s')$ tuple collected months ago can still contribute to today's gradient update.

On-policy methods discard data after each policy update. Once $\pi_\theta$ changes, old trajectories no longer provide unbiased gradient estimates, so the data must be thrown away.

\subsection{Engineering Efficiency Is a Different Story}

In practice, raw sample efficiency is only one factor:

\begin{center}
\begin{tabular}{@{} l p{5.5cm} p{5.5cm} @{}}
\toprule
\textbf{Factor} & \textbf{PPO (near on-policy)} & \textbf{SAC (off-policy)} \\
\midrule
Sample efficiency & Lower---data discarded after a few epochs & Higher---replay buffer reused extensively \\
Training stability & Very stable; robust to hyperparameters & Can diverge if Q-estimates blow up \\
Hyperparameter tuning & Easy; few sensitive knobs & Harder; temperature $\alpha$, replay ratio, critic architecture matter \\
Parallelism & Excellent; scales to thousands of workers & Moderate; replay buffer is a bottleneck \\
Reproducibility & High & Lower; sensitive to seeds and ordering \\
\bottomrule
\end{tabular}
\end{center}

\begin{keyidea}[Theoretical $\neq$ Engineering Efficiency]
Theoretical sample efficiency measures how many times each transition can be reused by the optimizer. Engineering efficiency measures whether, given the same engineering effort, you can stably obtain a good policy. PPO wastes samples but is easy to make work. SAC reuses samples but demands more careful tuning. If the off-policy critic diverges, reusing a million transitions is worthless.
\end{keyidea}

This is why PPO dominates in RLHF, large-scale simulation, and robotics despite its lower theoretical sample efficiency: stability and parallelizability often matter more than per-sample reuse.

%% ----------------------------------------------------------------
\section{The Compressed Decision Rule}
\label{sec:compressed-rule}

The entire on/off-policy distinction compresses to:

\begin{quote}
If the update formula requires the sample actions to be drawn from the current policy, the algorithm is \textbf{on-policy}. If it can use any historical actions as long as $(s, a, r, s')$ is available, it is \textbf{off-policy}. If it can use actions from a recent old policy but must constrain the divergence between old and new, it is \textbf{near on-policy}.
\end{quote}

The underlying question is always the same: \emph{is the action distribution in the historical data allowed to differ from the distribution of the policy currently being optimized?}
